%% kapa: exam=1 round=29 year=2018 period=2 form=A subject=appraisal_law subject_ko=감정평가관계법규 seq=24 number=24
%% source: Q-Net 감정평가사 기출문제 (cst00309) artlSeq=5202808 file=extracted/1008/5202808/2018년 제29회 감정평가사 1차시험 기출문제/2018년 제29회 감정평가사 1차시험 기출문제/2교시 A형.pdf
%% conversion: kordoc 4.24.51 -> tools/gen1.py (auto); review=proofread-sonnet; flags=none
\begin{question}{24}
건축법령상 건축 관련 입지와 규모의 사전결정에 관한 설명으로 옳지 않은 것은?
\begin{choices}
  \item 건축허가 대상 건축물을 건축하려는 자는 건축허가를 신청하기 전에 허가권자에게 해당 대지에 건축 가능한 건축물의 규모에 대한 사전결정을 신청할 수 있다.
  \item 사전결정신청자는 건축위원회 심의와「도시교통정비 촉진법」에 따른 교통영향평가서의 검토를 동시에 신청할 수 있다.
  \item 허가권자는 사전결정이 신청된 건축물의 대지면적이「환경영향평가법」에 따른 소규모 환경영향평가 대상사업인 경우 환경부장관이나 지방환경관서의 장과 소규모 환경영향평가에 관한 협의를 하여야 한다.
  \item 사전결정신청자가 사전결정 통지를 받은 경우에는「하천법」에 따른 하천점용허가를 받은 것으로 본다.
  \item 사전결정신청자는 사전결정을 통지받은 날부터 2년 이내에 건축허가를 받아야 하며, 이 기간에 건축허가를 받지 아니하면 사전결정의 효력은 상실된다.
\end{choices}
\end{question}
\answer{5}
